% ============================================================================
\chapter{Columnas de la hoja ENTRADA}\label{anx:columnas}
% ============================================================================

\renewcommand{\arraystretch}{1.25}
\begin{longtable}{@{}l>{\raggedright\arraybackslash}p{3.9cm}>{\raggedright\arraybackslash}p{4.9cm}>{\raggedright\arraybackslash}p{3.6cm}c@{}}
\toprule
\textbf{Col.} & \textbf{Nombre} & \textbf{Qué escribir} & \textbf{Documento fuente} & \textbf{Oblig.} \\ \midrule
\endhead
A & Nº & Nada: se numera solo & --- & --- \\
B & ID contenedor & Número del contenedor o del BL & BL / manifiesto & Sí \\
C & Fecha de arribo & dd/mm/aaaa & Aviso de arribo & Sí \\
D & Transitaria & ENCI, A.V u Otra (lista) & Documentos de la transitaria & Sí \\
E & Tipo & 20', 40', 40' HC u Otro (lista) & BL & No \\
F & Nº de bultos & Cantidad de bultos & Manifiesto & No \\
G & kg Occidente & kg con destino a Occidente & Desglose por destinatario & Sí \\
H & kg Centro & kg con destino a Centro & Desglose por destinatario & Sí \\
I & kg Oriente & kg con destino a Oriente & Desglose por destinatario & Sí \\
J & Viajes al puerto & Nº de viajes del camión propio & Hoja de ruta & Opc. \\
K & Litros reales al puerto & Total de litros & Vales de combustible & Opc. \\
L & Dieta pagada & Total en USD & Comprobante de dieta & Opc. \\
M & Jornales del desagrupe & Trabajadores $\times$ días & Nómina / control de asistencia & Opc. \\
N & Estadía / almacenaje & USD cobrados por demora & Factura de la naviera o de la transitaria & No \\
O & Litros distribución Occidente & Total de litros & Vales / hoja de ruta & Opc. \\
P & Litros traslado a Centro & Total de litros & Vales / hoja de ruta & Opc. \\
Q & Litros distribución Centro & Total de litros & Vales / hoja de ruta & Opc. \\
R & Jornales centro transitorio Centro & Trabajadores $\times$ días & Nómina / asistencia & Opc. \\
S & Litros traslado a Oriente & Total de litros & Vales / hoja de ruta & Opc. \\
T & Litros distribución Oriente & Total de litros & Vales / hoja de ruta & Opc. \\
U & Jornales centro transitorio Oriente & Trabajadores $\times$ días & Nómina / asistencia & Opc. \\
V & Otros gastos directos & USD & Factura / comprobante & No \\
W & Detalle de otros gastos & Texto breve & --- & No \\
X & Otros ingresos & USD & Factura emitida & No \\
Y & Fecha de fin de la distribución & dd/mm/aaaa & Última entrega & No \\
Z & Estado & Abierto o Cerrado (lista) & --- & Sí \\
AA & Observaciones & Números de documentos, explicaciones & --- & No \\
AB & Alertas & Nada: automático & --- & --- \\
AC & Utilidad antes de impuesto & Nada: automático & --- & --- \\
\bottomrule
\end{longtable}
\noindent\textbf{Opc.} = opcional: si se deja vacía, se usa la norma técnica de \hoja{PARAMETROS}.

% ============================================================================
\chapter{Glosario}\label{anx:glosario}
% ============================================================================

\begin{description}[style=nextline,leftmargin=1.2em,itemsep=4pt]
\item[Absorción (costeo por)] Repartir los gastos fijos entre los contenedores para conocer su costo completo.
\item[Alerta] Aviso automático de que falta un dato o de que algo se sale de lo normal.
\item[Celda] Cada casilla de una hoja, nombrada por su columna y su fila, como \celda{C5}.
\item[Centro de costo] Grupo al que se asigna un gasto: GENERAL, Occidente, Centro u Oriente.
\item[Conciliación] Comparación de dos registros que deberían coincidir: lo que declaran los gestores y lo que dicen las operaciones.
\item[Costo fijo] Gasto que se paga igual haya o no contenedores: alquiler, salarios de oficina, depreciación\dots
\item[Costo variable] Gasto que depende de cada contenedor: combustible, jornales, cita, comisión variable\dots
\item[Depreciación] Pérdida de valor del camión o de un equipo por el uso y el tiempo, repartida por meses.
\item[Desagrupe] Separar la carga del contenedor por destinatario.
\item[Fijos no absorbidos] Gastos fijos de un mes que ningún contenedor cubrió.
\item[Jornal] Un día de trabajo de una persona. 4 personas durante 3 días son 12 jornales.
\item[Margen de contribución] Ingresos menos costos variables.
\item[Margen de seguridad] Cuánto pueden bajar las ventas antes de tener pérdidas.
\item[Norma técnica] Consumo considerado normal (litros por viaje, kg por jornal\dots). Se usa cuando falta el dato real y para generar alertas.
\item[Período] Los 24 meses que abarca el libro, desde \texttt{Fecha\_Inicio}.
\item[Prorrateo] Reparto proporcional de un gasto, por ejemplo según los kg.
\item[Punto de equilibrio] Volumen con el que los ingresos cubren exactamente todos los gastos: ni se gana ni se pierde.
\item[Tarifa mínima de equilibrio] Precio por kg más bajo con el que una región no pierde dinero.
\item[UAI] Utilidad antes del impuesto sobre utilidades.
\item[Vigencia] Período durante el que rige un precio. Empieza en la fecha de su fila en \hoja{PRECIOS}.
\end{description}

% ============================================================================
\chapter{Lista de comprobación mensual (para imprimir)}\label{anx:checklist}
% ============================================================================

\noindent Mes que se cierra: \rule{5cm}{0.4pt} \hfill Fecha del cierre: \rule{4cm}{0.4pt}

\bigskip
\renewcommand{\arraystretch}{2.0}
\begin{tabularx}{0.96\linewidth}{@{}c>{\raggedright\arraybackslash}X p{3.4cm}@{}}
\toprule
& \textbf{Tarea} & \textbf{Observaciones} \\ \midrule
$\square$ & Todos los contenedores del mes registrados y en «Cerrado» (\hoja{ENTRADA}) & \\
$\square$ & Alertas resueltas o explicadas & \\
$\square$ & Gastos fijos reales del mes escritos (\hoja{COSTOS\_FIJOS}) & \\
$\square$ & Cambios de precio registrados como filas nuevas (\hoja{PRECIOS}) & \\
$\square$ & Base variable de los 21 gestores escrita (\hoja{COMERCIAL}) & \\
$\square$ & Conciliación de comercialización: CUADRA o diferencia explicada & \\
$\square$ & Cuadres de \hoja{INICIO}: OK y OK & \\
$\square$ & Resultado revisado: utilidad, costo/kg, equilibrio, regiones & \\
$\square$ & Informe del \hoja{DASHBOARD} exportado a PDF & \\
$\square$ & Archivo guardado y copia de respaldo hecha (USB, correo o nube) & \\
\bottomrule
\end{tabularx}

\vspace{1cm}
\noindent Utilidad antes de impuesto del mes: \rule{4cm}{0.4pt}\hfill Contenedores del mes: \rule{3cm}{0.4pt}

\vspace{14pt}
\noindent Firma: \rule{5cm}{0.4pt}

% ============================================================================
\chapter{Cómo actualizar esta guía}\label{anx:mantener}
% ============================================================================

Esta guía está escrita en LaTeX para que cualquier cambio sea sencillo y la guía siga siempre coherente con el libro de Excel.

\section{Dónde está cada cosa}
\begin{center}
\begin{tabularx}{\linewidth}{@{}l X@{}}
\toprule
\textbf{Archivo} & \textbf{Contenido} \\ \midrule
\texttt{guia/guia\_usuario.tex} & Archivo principal: portada, orden de los capítulos y versión (\verb|\version|). \\
\texttt{guia/capitulos/*.tex} & Un archivo por capítulo. Para cambiar un texto, edite el capítulo correspondiente. \\
\texttt{guia/estilo.sty} & Colores, cuadros (nota, atención, consejo, resultado) y comandos (\verb|\hoja|, \verb|\celda|, \verb|\escriba|\dots). \\
\texttt{guia/valores.tex} & Números de los ejercicios. \textbf{Se genera solo}: no lo edite a mano. \\
\texttt{guia/img/} & Capturas del libro. \textbf{Se generan solas}. \\
\texttt{build/capturas.py} & Define los ejercicios, las capturas (hoja, rango, recuadros) y los valores citados. \\
\bottomrule
\end{tabularx}
\end{center}

\section{Volver a generar la guía}
Requiere Python con \texttt{openpyxl}, \texttt{lxml} y \texttt{Pillow}, LibreOffice Calc, \texttt{pdftoppm} y una distribución de LaTeX (TeX Live o MiKTeX).
\begin{verbatim}
build/construir.sh          # 1. regenera el libro de Excel
build/construir_guia.sh     # 2. capturas + valores + PDF
\end{verbatim}
Si cambia el libro de Excel (una hoja, una columna, una fórmula), ejecute ambos pasos: las imágenes y los números de la guía se actualizan solos.
Si un cambio mueve una hoja o una columna, ajuste también los rangos de \texttt{build/capturas.py} y las celdas que nombra el texto.

\section{Comandos de la guía}
\begin{center}
\begin{tabular}{@{}p{7.2cm}p{8.4cm}@{}}
\toprule
\textbf{Escriba en el .tex} & \textbf{Resultado} \\ \midrule
\verb|\hoja{ENTRADA}| & \hoja{ENTRADA} \\
\verb|\celda{C5}| & \celda{C5} \\
\verb|\escriba{5000}| & \escriba{5000} \\
\verb|\tecla{Enter}| & \tecla{Enter} \\
\verb|\recuadro{1}| & \recuadro{1} \\
\verb|\captura{ej2_entrada}{Pie}| & Inserta \texttt{img/ej2\_entrada.png} con su pie de figura \\
\verb|\begin{resultado}...\end{resultado}| & Cuadro verde «Lo que debe ver si lo hizo bien» \\
\bottomrule
\end{tabular}
\end{center}
